% ANCHOR: end_of_section 0.1
\begin{remark}[Physics notation: fields, action and propagator]
In the components of the physics notation key after Section 1.2:
\begin{itemize}
\item $x^m$, $m=1,\dots,4$, are the coordinates of Euclidean $\mathbb R^4$; $W_{\mathbb R}^\perp$ carries the $SO(6)_R$ vector index $I=1,\dots,6$ and $W$, the $\mathbf 4$ of $SU(4)_R$, the index $i=1,\dots,4$.
\item $A|_V=A_m$ is the gauge field, $\iota_wA=w_IA_I$ the six scalars, $\lambda=(\lambda_\alpha{}^i,\bar\lambda_{\dot\alpha i})$ the four gauginos, $c$ the Faddeev--Popov ghost, and $A^*_m$, $A^*_I$, $\lambda^{*\alpha}{}_i$, $\bar\lambda^{*\dot\alpha i}$, $c^*$ the antifields, all $\mathfrak g$-valued.
\item $S$ is the BV action of Euclidean $\mathcal N=4$ super Yang--Mills theory with weight $e^{S/\hbar}$, $Q$ its linearised BRST operator, $I$ the interaction and $\{-,-\}$ the BV antibracket.
\item $Q^{GF}$ is the Feynman-gauge gauge-fixing operator, $H=-\partial_m\partial_m$ ($|p|^2$ in momentum space), $K_t$ the heat kernel with momentum factor $e^{-t|p|^2}$, and $P(\varepsilon,\infty)$ the Feynman-gauge propagator with proper-time cutoff $\varepsilon$ (mass dimension $-2$): each component is its fibre factor times $\int_\varepsilon^\infty e^{-t|p|^2}dt=e^{-\varepsilon|p|^2}/|p|^2$. For $\mathfrak{su}(N)$, with $A_M=-igA^a_MT^a$, $M=(m,I)$, and $\lambda=-\tfrac{ig}{\sqrt2}\lambda^aT^a$, the edge $\hbar P(\varepsilon,\infty)$ of $\hbar\partial_{P(\varepsilon,\infty)}$ has the components
\[
A^a_M\otimes A^b_N:\ \delta^{ab}\delta_{MN}\,\frac{e^{-\varepsilon|p|^2}}{|p|^2},\qquad \lambda^a_\alpha{}^i\otimes\bar\lambda^b_{\dot\beta j}:\ -\delta^{ab}p_m\sigma^m_{\alpha\dot\beta}\delta^i_j\,\frac{e^{-\varepsilon|p|^2}}{|p|^2}.
\]
\item $W(P,J)$ is $\hbar$ times the sum of connected Feynman graphs with edges $\hbar P$ and vertices $J/\hbar$, each graph weighted by $\hbar$ to its number of loops; $\varepsilon^{-1}$ and $\log\varepsilon$ are the quadratic and logarithmic divergences.
\item With $\omega=f\,d^4x$, $\int\omega\,(\iota_{w_1}A,\iota_{w_2}A)=w_{1I}w_{2J}\int d^4x\,f\,(A_I,A_J)$, equal to $2\hbar\,w_{1I}w_{2J}\int d^4x\,f\,A^a_IA^a_J$ for $\mathfrak{su}(N)$; $\delta$ is its source, $O(\varepsilon)$ the bare operator including its counterterm, and $O[\infty]$ the renormalised operator.
\item $\Delta=2E=2x^m\partial_m$, and $[(d_{CE}+Q)O[\infty]](x_1)$ is the quantum-corrected transformation of $O$ along the generator $x_1$ of the key: $P_m$, $K_m$, $\Delta$, rotations, R-symmetry, $Q^i_\alpha$, $\bar Q_{\dot\alpha i}$, $S^\alpha{}_i$, $\bar S^{\dot\alpha i}$.
\end{itemize}
\end{remark}

% ANCHOR: end_of_section 0.2
\begin{remark}[Physics notation: the UV boundary]
\begin{itemize}
\item $K_0$ is $-\delta^4(x-y)$ times the inverse fibre pairing on each field--antifield pair; $K_t$ has $\delta^4(x-y)$ replaced by $(4\pi t)^{-2}e^{-|x-y|^2/4t}$, momentum factor $e^{-t|p|^2}$, so $K_\infty=0$ and $K_\varepsilon-K_0$ carries the factor $e^{-\varepsilon|p|^2}-1$.
\item $L_t=(Q^{GF}\otimes1)K_t$ is the integrand of the proper-time representation $P(\varepsilon,\infty)=\int_\varepsilon^\infty L_t\,dt$: the fibre factor of the propagator times $e^{-t|p|^2}$.
\item $P(0,\infty)$ is the Feynman-gauge propagator without cutoff, fibre factor times $1/|p|^2$, and $Q\,P(0,\infty)\star+P(0,\infty)\star\,Q=1$ makes $P(0,\infty)\star$ a Green's operator of the free BRST operator $Q$.
\item $\Delta_\varepsilon=-\partial_{K_\varepsilon}$ is the BV Laplacian regularised by the heat kernel at proper time $\varepsilon$; the quantum master equation at the UV boundary reads $(d_{CE}+Q)J+\tfrac12\{J,J\}_\varepsilon+\hbar\Delta_\varepsilon J=0$, with classical part the Zinn-Justin equation.
\item The $\varepsilon^{-1}$ mass terms are $(A_m,A_m)+(A_I,A_I)$, and the finite counterterm forced by the master equation is $(\partial_mA_m,\partial_nA_n)$ (Section 0.4).
\end{itemize}
\end{remark}

% ANCHOR: end_of_section 0.3
\begin{remark}[Physics notation: the objects of the chapters]
The Feynman-gauge operator of Chapter 3 (Definition \ref{def:c3-qgf}) is, with the slot of the image in brackets,
\[
\begin{aligned}
&A^*_M\mapsto-2A^*_M\ [A_M],\qquad A_m\mapsto\partial_mA_m\ [c],\qquad c^*\mapsto\partial_mc^*\ [A^*_m],\\
&\lambda^{*\alpha}{}_i\mapsto-i\,\partial_m\lambda^{*\alpha}{}_i\sigma^m_{\alpha\dot\alpha}\ [\bar\lambda_{\dot\alpha i}],\qquad \bar\lambda^{*\dot\alpha i}\mapsto i\,\sigma^m_{\alpha\dot\alpha}\partial_m\bar\lambda^{*\dot\alpha i}\ [\lambda_\alpha{}^i],
\end{aligned}
\]
with $H=[Q,Q^{GF}]=-\partial_m\partial_m$. The single letters are $\int\omega\,(X,A_I)$, $\int\omega\,(X,v_mA_m)$, $\int\omega\,(\zeta,\lambda)$, $\int\omega\,(X,c)$; the scalar bilinears are $w_{1I}w_{2J}\int\omega\,(A_I,A_J)$, whose traceless part is the $\mathbf{20'}$ of $SU(4)_R$ and whose trace $\int\omega\,(A_I,A_I)$ is the Konishi operator. The other superconformal components of Chapter 7 are those along $Q^i_\alpha$, $\bar Q_{\dot\alpha i}$, $S^\alpha{}_i$, $\bar S^{\dot\alpha i}$ and $K_m$.
\end{remark}

% ANCHOR: end_of_section 0.4
\begin{remark}[Physics notation: the results of Step 1]
With $(A_I,A_I)=\tfrac14(A^{ij},A_{ij})$, $F^0_{mn}=\partial_mA_n-\partial_nA_m$ and $d^*A|_V=-\partial_mA_m$, the quadratic counterterm of Chapter 4 is
\[
\begin{aligned}
\frac{\hbar\operatorname{tr}_{\mathfrak g}(\mathrm{ad}_X\mathrm{ad}_Y)}{16\pi^2(X,Y)}\int d^4x\Big\{&\frac9{4\varepsilon}\big[(A_m,A_m)+(A_I,A_I)\big]\\
&+\log\varepsilon\Big[-\tfrac12(F^0_{mn},F^0_{mn})-(\partial_mA_I,\partial_mA_I)+8i\big(\bar\lambda_{\dot\alpha i},\bar\sigma^{m\dot\alpha\alpha}\partial_m\lambda_\alpha{}^i\big)-(\partial_mc,A^*_m)\Big]\Big\},
\end{aligned}
\]
gauge-boson and scalar mass terms at order $\varepsilon^{-1}$ and kinetic terms at order $\log\varepsilon$. The $\log\varepsilon$ part through quartic order is the unit times $\log\varepsilon\,(Q+\{I,-\})$ of
\[
\int d^4x\Big\{2(A_m,A^*_m)+2(A_I,A^*_I)-4\big[(\lambda_\alpha{}^i,\lambda^{*\alpha}{}_i)-(\bar\lambda_{\dot\alpha i},\bar\lambda^{*\dot\alpha i})\big]-(c^*,c)\Big\},
\]
the BV form of a wave-function renormalisation, with vanishing one-loop beta function of $g$. The finite counterterm is $-\tfrac{27}{16}$ units times $\int d^4x\,(\partial_mA_m,\partial_nA_n)$, removing the obstruction $-\tfrac{27}{8}$ units times $\int d^4x\,(\partial^2c,\partial_mA_m)$. In Chapters 5 and 6, $(\zeta,\lambda)=(\lambda_\alpha{}^i,\zeta^\alpha{}_i)-(\bar\lambda_{\dot\alpha i},\bar\zeta^{\dot\alpha i})$, $(w_1,w_2)=w_{1I}w_{2I}$, the traceless bilinear is the $\mathbf{20'}$ and the trace $\int\omega\,(A_I,A_I)=-\tfrac{4\hbar}{g^2}\int\omega\operatorname{tr}(A_IA_I)$ the Konishi operator. The coefficient times the unit is the one-loop anomalous dimension, and $w_O$ plus it the scaling dimension; for $\mathfrak{su}(N)$ the table gives the anomalous dimensions $0$ for the $\mathbf{20'}$, $\tfrac{3g^2N}{4\pi^2}$ for the Konishi operator and, in Feynman gauge, $\tfrac{g^2N}{8\pi^2}$ for $A_I$ and $A_m$, $\tfrac{g^2N}{4\pi^2}$ for $\lambda$, $\tfrac{g^2N}{16\pi^2}$ for $c$.
\end{remark}
